% Irodov, user-supplied Problems in General Physics, printed pp.33--34
% and continuations on printed pp.34--35, source PDF pp.42--44.
% In this chapter G is the gravitational constant and M_s is the solar mass.
\exproblem{200}{33}
一颗行星在太阳引力作用下作圆周运动，日心系速率为
$v=34.9\,\mathrm{km/s}$。求该行星的公转周期。
\exsolution
记太阳质量 $M_\odot$，引力参数
$\mu=GM_\odot$。圆轨道上的引力提供向心加速度：
\[
\frac{v^2}{r}=\frac{\mu}{r^2}
\Longrightarrow r=\frac{\mu}{v^2}.
\]
公转周期 $T=2\pi r/v$，故
\[
\boxed{T=\frac{2\pi GM_\odot}{v^3}}.
\]
取 $\mu\simeq1.3271\times10^{20}\,\mathrm{m^3/s^2}$、
$v=3.49\times10^4\,\mathrm{m/s}$，
算得 $T\simeq2.27\times10^2\,\mathrm d$（约 $227$ 天）。
\excheck 原书简答给出约 $225$ 天，属于所取太阳常数与
有效数字的近似差别；周期计算必须有 $v^3$ 在分母，
否则量纲不为时间。

\exproblem{201}{33}
木星绕太阳的周期是地球的 $12$ 倍，视二者轨道为圆。
求：（a）木星日心距离是地球的几倍；
（b）木星在日心系中的速率、向心加速度。
\exsolution
开普勒第三定律 $T^2\propto r^3$，
故
\[
\boxed{\frac{r_J}{r_E}=12^{2/3}\simeq5.24}.
\]
由圆轨道速率 $v=2\pi r/T$ 与
$a=v^2/r$，
\[
\frac{v_J}{v_E}=12^{-1/3},\qquad
\frac{a_J}{a_E}=12^{-4/3}.
\]
取地球公转速率 $v_E\simeq29.8\,\mathrm{km/s}$、
向心加速度 $a_E\simeq5.93\times10^{-3}\,\mathrm{m/s^2}$：
\[
\boxed{v_J\simeq13.0\,\mathrm{km/s},\qquad
a_J\simeq2.16\times10^{-4}\,\mathrm{m/s^2}}.
\]
\excheck 木星轨道半径较大、速率和向心加速度均较小；
这里用的是日心近圆轨道近似，不是木星相对地球的速度。

\exproblem{202}{33}
行星绕太阳沿椭圆轨道运动，
其最近、最远日心距离分别为 $r$ 和 $R$。
利用开普勒定律求公转周期。
\exsolution
椭圆的半长轴
$a=(r+R)/2$。
按开普勒第三定律，与同一太阳引力参数下
半径为 $a$ 的圆轨道具有同样的周期：
\[
T^2=\frac{4\pi^2a^3}{GM_\odot}.
\]
因此
\[
\boxed{T=2\pi\sqrt{\frac{[(r+R)/2]^3}{GM_\odot}}
=\pi\sqrt{\frac{(r+R)^3}{2GM_\odot}}}.
\]
\excheck 当 $R=r$，表达式退化为
$T=2\pi\sqrt{r^3/(GM_\odot)}$，
恰为圆轨道公式。

\exproblem{203}{33--34}
一小天体在日心参考系内，从与地球公转轨道半径相等
的距离 $R_E$ 处由静止开始落向太阳。
利用开普勒定律估算下落时间，忽略太阳表面半径。
\exsolution
径直下落可看作离心率趋近 $1$ 的退化椭圆轨道的半程：
远日点为 $R_E$，近日点趋于 $0$，半长轴为 $R_E/2$。
其完整退化椭圆的形式周期与同半长轴的圆轨道相同：
\[
T_*=\frac{T_E}{2^{3/2}},
\]
其中 $T_E$ 是地球一年公转周期。
从远日点落到近日点只经过半个周期，故
\[
\boxed{t_{\rm fall}=\frac{T_E}{4\sqrt2}
\simeq64.6\,\mathrm d\approx65\,\mathrm d}.
\]
\excheck 忽略了太阳有限半径；若以太阳表面为终点，
实际下落时间会略短于该中心点极限。

\exproblem{204}{34}
把太阳系所有线性尺寸按同一比例 $\eta$ 缩小，
所用材料的平均密度仍与原恒星、行星各自的平均密度相同。
问行星模型公转周期相对原系统如何变化？
\exsolution
设所有长度缩小为原来的 $1/\eta$
（若 $\eta>1$）。平均密度不变意味着各质量按
体积缩放，故中心天体质量
$M'_\odot=M_\odot/\eta^3$。
行星轨道半长轴 $a'=a/\eta$，
开普勒第三定律给出
\[
\frac{{T'}^2}{T^2}
=\frac{(a')^3/M'_\odot}{a^3/M_\odot}
=\frac{(a^3/\eta^3)/(M_\odot/\eta^3)}
{a^3/M_\odot}=1.
\]
因此
\[
\boxed{T'=T}.
\]
\excheck 系统公转速率按线性尺寸同比缩小，
但周期在相同平均密度的相似模型中保持不变。

\exproblem{205}{34}
一对双星在万有引力下彼此绕共同质心作圆周运动。
已知两星总质量为 $M$、公转周期为 $T$，
求两星之间的距离 $\ell$。
\exsolution
设两星质量为 $m_1,m_2$，间距为 $\ell$，
共同角速度 $\omega=2\pi/T$。
关于质心的轨道半径分别为
$r_1=m_2\ell/M$、$r_2=m_1\ell/M$。
对第一星列向心力方程：
\[
\frac{Gm_1m_2}{\ell^2}
=m_1\omega^2\frac{m_2}{M}\ell.
\]
约去质量因子得
$\omega^2=GM/\ell^3$，所以
\[
\boxed{\ell=\left(\frac{GMT^2}{4\pi^2}\right)^{1/3}}.
\]
\excheck $M=m_1+m_2$ 是两星\emph{总质量}，
并非任意一颗星的质量；结论与质量比无关。

\exproblem{206}{34}
求引力相互作用势能：（a）两个质量分别为 $m_1,m_2$、
间距 $r$ 的质点；（b）质量为 $m$ 的质点位于
一均匀细杆轴线的延长线上，距杆近端为 $a$，
细杆质量 $M$、长度 $\ell$。
再求（b）的引力大小。
\exsolution
取无穷远相互作用势能为零。
（a）按牛顿万有引力积分：
\[
\boxed{U_a(r)=-\frac{Gm_1m_2}{r}}.
\]
（b）均匀细杆线密度 $\lambda=M/\ell$，
从杆近端起坐标 $x\in[0,\ell]$，
微元质量为 $\mathrm dm=\lambda\,\mathrm dx$。
质点与微元距离为 $a+x$，故
\[
U_b=-Gm\int_0^\ell\frac{\lambda\,\mathrm dx}{a+x}
=\boxed{-\frac{GmM}{\ell}
\ln\left(1+\frac{\ell}{a}\right)}.
\]
朝向细杆的引力大小由势能对间隔的导数得到：
\[
\boxed{F=Gm\int_0^\ell
\frac{\lambda\,\mathrm dx}{(a+x)^2}
=\frac{GmM}{a(a+\ell)}}.
\]
\excheck $\ell\to0$ 而总质量 $M$ 不变时，
$U_b\to-GmM/a$，$F\to GmM/a^2$，
退化为两个质点。

\exproblem{207}{34}
一质量为 $m$ 的行星绕太阳作椭圆轨道，
最远、最近日心距离分别为 $r_1,r_2$。
求相对太阳中心的轨道角动量大小。
\exsolution
记太阳引力参数 $\mu=GM_\odot$，
近、远日点切向速率分别为 $v_2,v_1$。
角动量守恒：
$L=mr_1v_1=mr_2v_2$。
机械能守恒：
\[
\frac{L^2}{2mr_1^2}-\frac{GmM_\odot}{r_1}
=\frac{L^2}{2mr_2^2}-\frac{GmM_\odot}{r_2}.
\]
移项并约分可得
$L^2/(2m)=GmM_\odot r_1r_2/(r_1+r_2)$，
因此
\[
\boxed{L=m\sqrt{\frac{2GM_\odot r_1r_2}{r_1+r_2}}}.
\]
\excheck 若 $r_1=r_2=r$，公式化为
$L=m\sqrt{GM_\odot r}$，正是圆轨道角动量。

\exproblem{208}{34}
利用守恒定律证明：质量为 $m$ 的行星绕太阳作椭圆轨道时，
其总机械能仅由椭圆半长轴 $a$ 决定；求 $E(a)$。
\exsolution
仍记 $\mu=GM_\odot$，最近、最远距离
分别为 $r_p,r_a$，所以 $a=(r_p+r_a)/2$。
在两个拱点径向速度均为零，角动量 $L$
及总机械能相同：
\[
E=\frac{L^2}{2mr_p^2}-\frac{\mu m}{r_p}
=\frac{L^2}{2mr_a^2}-\frac{\mu m}{r_a}.
\]
先消去 $E$ 可得
$L^2/(2m)=\mu m r_pr_a/(r_p+r_a)$；
代入任一拱点：
\[
E=\frac{\mu m r_a}{r_p(r_p+r_a)}
-\frac{\mu m}{r_p}
=-\frac{\mu m}{r_p+r_a}.
\]
故
\[
\boxed{E=-\frac{GM_\odot m}{2a}}.
\]
\excheck $E$ 为负表示引力束缚轨道，
且只依赖半长轴、不单独依赖偏心率。

\exproblem{209}{34}
行星在椭圆轨道上的某一时刻距离太阳为 $r_0$、
速率为 $v_0$，速度与日心径矢夹角为 $\alpha$。
求这一轨道的最近和最远日心距离。
\exsolution
记 $\mu=GM_\odot$，单位质量的守恒能量
$\varepsilon=v_0^2/2-\mu/r_0<0$，
比角动量 $h=r_0v_0\sin\alpha$。
定义无量纲参数
$q=r_0v_0^2/\mu$，椭圆条件 $0<q<2$。
半长轴由 $\varepsilon=-\mu/(2a)$ 得
$a=r_0/(2-q)$。
偏心率由能量—角动量关系
\[
e^2=1+\frac{2\varepsilon h^2}{\mu^2}
=1-q(2-q)\sin^2\alpha.
\]
所以两个拱点距离为
\[
\boxed{r_{\min,\max}=
\frac{r_0}{2-q}
\left[1\mp\sqrt{1-q(2-q)\sin^2\alpha}\right]}.
\]
\excheck 负号对应近日点，正号对应远日点；
若输入参数不满足束缚轨道条件，不可称远日点有限。

\exproblem{210}{34}
一宇宙天体在距太阳很远处以速率 $v_0$
向太阳接近，其入射渐近直线距太阳中心的冲量参数
（瞄准距离）为 $\ell$（图1.51）。求最近日心距离。
\exsolution
记 $\mu=GM_\odot$，质点质量在比能量与比角动量中相消。
在无穷远处有
$\varepsilon=v_0^2/2$、$h=\ell v_0$。
最近点处径向速度为零，令该距离为 $r_{\min}$。
由能量守恒
\[
\frac{v_0^2}{2}
=\frac{h^2}{2r_{\min}^2}-\frac{\mu}{r_{\min}},
\]
于是
$v_0^2r_{\min}^2+2\mu r_{\min}-\ell^2v_0^2=0$。
取正根：
\[
\boxed{r_{\min}=\frac{\mu}{v_0^2}
\left[\sqrt{1+\left(\frac{\ell v_0^2}{\mu}\right)^2}-1\right]}.
\]
\excheck 当 $\ell=0$，理想质点可直冲中心；
当 $\mu\to0$ 时应以有理化形式取极限，
可得 $r_{\min}\to\ell$，符合自由直线运动。
